% ============================================================
% Prueba Técnica — Resultados
% Choque petrolero de 2022 y subsidios a combustibles fósiles
% en América Latina y el Caribe
% Candidato: Daniel Mendivelso
% ------------------------------------------------------------
% Compilar:  xelatex -interaction=nonstopmode Prueba_Tecnica_Resultados.tex
% Para ver las NOTAS DEL EXPOSITOR, descomentar la línea:
%     \setbeameroption{show notes}
% (o "show notes on second screen" para presentar con notas).
% ============================================================
\documentclass[aspectratio=169,10pt]{beamer}

% --- Tipografía Unicode (XeLaTeX) ---
% IMPORTANTE: compilar con xelatex. fontspec maneja bien los signos
% ¿ ¡ á é í ... (con inputenc/fontenc el ¿ se renderiza como £).
\usepackage{fontspec}
\setsansfont{texgyreheros}[              % clon de Helvetica, working-paper limpio
  Extension      = .otf,
  UprightFont    = *-regular,
  ItalicFont     = *-italic,
  BoldFont       = *-bold,
  BoldItalicFont = *-bolditalic,
  Scale          = 0.92]
\renewcommand{\familydefault}{\sfdefault}
\usepackage[spanish]{babel}
\usepackage{microtype}

% --- Paquetes ---
\usepackage{amsmath}
\usepackage{graphicx}
\usepackage{booktabs}
\usepackage{xcolor}

% ============================================================
% PALETA "Gris carbón minimal" (sin marca institucional)
% ============================================================
\definecolor{Carbon}{HTML}{333333}     % acento principal
\definecolor{GrisMedio}{HTML}{6E6E6E}   % detalle / secundario
\definecolor{AcentoRojo}{HTML}{B00020}  % UNA cifra clave por slide
\definecolor{GrisClaro}{HTML}{EDEDED}   % fondos sutiles

% --- Tema limpio construido a mano ---
\usetheme{default}
\setbeamertemplate{navigation symbols}{}
\setbeamertemplate{itemize items}[circle]
\setbeamertemplate{itemize subitem}[triangle]

% Colores del tema
\setbeamercolor{structure}{fg=Carbon}
\setbeamercolor{frametitle}{fg=Carbon}
\setbeamercolor{title}{fg=Carbon}
\setbeamercolor{subtitle}{fg=GrisMedio}
\setbeamercolor{author}{fg=GrisMedio}
\setbeamercolor{date}{fg=GrisMedio}
\setbeamercolor{normal text}{fg=Carbon}
\setbeamercolor{itemize item}{fg=Carbon}
\setbeamercolor{itemize subitem}{fg=GrisMedio}
% Énfasis (\alert) en carbón negrita; el rojo se reserva para \key (UNA cifra por slide)
\setbeamercolor{alerted text}{fg=Carbon}
\setbeamerfont{alerted text}{series=\bfseries}
\setbeamercolor{block title}{bg=GrisClaro,fg=Carbon}
\setbeamercolor{block body}{bg=GrisClaro!40,fg=Carbon}

\setbeamerfont{frametitle}{series=\bfseries}
\setbeamerfont{title}{series=\bfseries,size=\LARGE}
\setbeamerfont{block title}{series=\bfseries}

% --- Frametitle: título en carbón + regla fina roja debajo ---
\setbeamertemplate{frametitle}{%
  \vspace{0.6em}%
  \begin{beamercolorbox}[wd=\paperwidth,leftskip=0.6cm,rightskip=0.6cm]{frametitle}%
    {\bfseries\insertframetitle}\par%
    \vspace{0.15em}%
    {\color{AcentoRojo}\rule{2.2cm}{1.2pt}}%
  \end{beamercolorbox}%
  \vspace{0.2em}%
}

% --- Footline minimalista en gris ---
\setbeamertemplate{footline}{%
  \hbox{%
    \begin{beamercolorbox}[wd=\paperwidth,ht=2.2ex,dp=1ex,leftskip=0.6cm,rightskip=0.6cm]{}%
      \usebeamerfont{date}\color{GrisMedio}\scriptsize
      Choque petrolero 2022 y subsidios fósiles en ALC%
      \hfill Daniel Mendivelso%
      \hfill \insertframenumber%
    \end{beamercolorbox}%
  }%
}

% --- Comando para resaltar la cifra clave (rojo, negrita) ---
\newcommand{\key}[1]{{\bfseries\color{AcentoRojo}#1}}

% --- Atajos de ruta a los assets ---
% figures/cropped/  = figuras 8-bit RGB con la banda de notas recortada
%                     (la nota va como texto LaTeX bajo cada figura).
% figures/compiled/ = figuras 8-bit RGB completas (con su nota quemada),
%                     usadas en los slides de respaldo.
% Las originales de 16-bit no renderizan bajo XeLaTeX/xdvipdfmx.
\newcommand{\fig}[1]{figures/cropped/#1}
\newcommand{\figfull}[1]{figures/compiled/#1}
\newcommand{\tab}[1]{tables/#1}

% Nota al pie de figura: gris, pequeña, JUSTIFICADA A LA IZQUIERDA,
% interlineado compacto. Va dentro de la columna (hereda su ancho).
\newcommand{\fignota}[1]{\par\vspace{0.2em}%
  {\scriptsize\color{GrisMedio}\setlength{\baselineskip}{0.85\baselineskip}%
   \raggedright #1\par}}

% Igual, pero para usar dentro de un \begin{center} (figura full-width):
% envuelve la nota en un minipage del ancho del texto, alineada a la izquierda.
\newcommand{\fignotaC}[1]{\par\vspace{0.2em}%
  \begin{minipage}{0.92\linewidth}%
    {\scriptsize\color{GrisMedio}\setlength{\baselineskip}{0.85\baselineskip}%
     \raggedright #1\par}%
  \end{minipage}}

% Caption tipo paper para figuras/tablas: "Figura N. Título descriptivo".
% Va ENCIMA del \includegraphics. #1 = etiqueta (Figura 1 / Tabla 5), #2 = título breve.
% La numeración sigue el nombre de archivo (fig1..5, tab1,2,4,5,6), no el orden de aparición.
\newcommand{\captitulo}[2]{%
  {\footnotesize\color{Carbon}\textbf{#1.}~#2\par}\vspace{0.2em}}

% --- Divisor de sección (reusa el patrón de "Material de respaldo") ---
% #1 = número de bloque, #2 = título, #3 = subtítulo. Frame [plain] sin footline.
\newcommand{\divisor}[3]{%
  {\setbeamertemplate{footline}{}%
   \begin{frame}[plain]
     \vfill
     \begin{center}
       {\color{AcentoRojo}\rule{2.6cm}{1.2pt}}\\[0.7em]
       {\scriptsize\color{GrisMedio}Bloque #1 \,/\, 7}\\[0.6em]
       {\Large\color{Carbon}\bfseries #2}\\[0.45em]
       {\color{GrisMedio}#3}\\[0.7em]
       {\color{AcentoRojo}\rule{2.6cm}{1.2pt}}
     \end{center}
     \vfill
   \end{frame}}}

% --- Espaciado de listas más aireado ---
\setlength{\parskip}{2pt}

% ============================================================
% Para PRESENTAR CON NOTAS, descomentar una de estas líneas:
% \setbeameroption{show notes}
% \setbeameroption{show notes on second screen=right}
% ============================================================

% --- Metadatos ---
\title{El choque petrolero de 2022 y los subsidios a combustibles fósiles}
\subtitle{Un experimento natural en América Latina y el Caribe}
\author{Daniel Mendivelso}
\date{2026}

\begin{document}

% ============================================================
% PORTADA
% ============================================================
{
\setbeamertemplate{footline}{} % portada sin footline
\begin{frame}[plain]
  \vfill
  \begin{center}
    {\color{AcentoRojo}\rule{3.5cm}{1.4pt}}\\[1.1em]
    {\usebeamerfont{title}\color{Carbon}\bfseries El choque petrolero de 2022\\[0.2em]
     y los subsidios a combustibles fósiles}\\[2.4em]
    {\large\color{Carbon}\textbf{Daniel Mendivelso}}\\[0.4em]
    {\color{GrisMedio}2026}\\[1.1em]
    {\color{AcentoRojo}\rule{3.5cm}{1.4pt}}
  \end{center}
  \vfill
\end{frame}
}

% ############################################################
% BLOQUE 1 — EL PROBLEMA
% ############################################################
\divisor{1}{El choque de 2022 y la pregunta de política}{El problema}

% --- Slide 1 (existente): El choque de 2022 ---
\begin{frame}{El choque de 2022: un experimento natural sobre los subsidios}
  \begin{itemize}
    \item En 2022 el Brent saltó de \key{USD 71 a 101 por barril (+42\%)} en promedio anual:
          el mayor salto de la serie 2015--2023.
    \item Detonante: la \alert{invasión de Rusia a Ucrania} (feb.\ 2022). Un choque
          \textbf{global, inesperado} y ajeno a la política de subsidios de cualquier país.
    \item Es \textbf{plausiblemente exógeno} a la política de subsidios de la región: ahí está
          la oportunidad de identificación.
  \end{itemize}

  \begin{block}{Pregunta y tesis}
    \textbf{¿Cómo respondieron los subsidios fósiles de la región y qué implicaciones fiscales
    tiene esa respuesta?}\\[0.3em]
    \textbf{Tesis:} el choque petrolero de 2022 no afectó a todos los países de la región por igual;
    su impacto fiscal, expresado a través del subsidio explícito a los combustibles, dependió de la
    \textbf{posición de exposición petrolera neta} de cada país.
  \end{block}

  \note{%
    El precio del Brent lo fija el mercado mundial; ningún país de la región lo mueve con su
    política de subsidios. Eso es lo que lo hace exógeno y útil como fuente de variación. La
    pregunta no es solo ``¿subieron los subsidios?'' sino ``¿quién y por qué?'', porque ahí
    está la lección de política. Anticipo la tesis: la exposición petrolera neta del país
    determina la dirección del golpe, y esa es la variable sobre la que construyo todo.}
\end{frame}

% --- Slide 1b (NUEVO): experimento natural + describir vs medir + contexto FMI ---
\begin{frame}{Por qué este episodio permite medir un efecto, no solo describirlo}
  \small
  \begin{block}{El choque como experimento natural}
    El precio internacional lo fija el mercado mundial; \textbf{ningún país de la región lo mueve
    con su política de subsidios}. Eso \alert{rompe la causalidad inversa}: la variación del
    tratamiento no viene de las decisiones que queremos estudiar, sino de afuera.
  \end{block}

  \textbf{Describir $\neq$ medir.} Ver que el subsidio subió en 2022 \emph{describe} una
  evolución. Pero la pregunta es \textbf{causal} (¿cuánto se debe al choque?), y una pregunta
  causal exige un \textbf{contrafactual}: qué habría pasado sin el choque. Para eso necesito un
  grupo de comparación, no solo una serie en el tiempo.\\[0.3em]
  \textbf{Tesis que guía todo:} el impacto fiscal del choque, medido por el subsidio explícito,
  dependió de la \key{posición de exposición petrolera neta} de cada país. Es lo que separa a quien
  gana renta de quien solo sufre el costo.

  \begin{block}{Contexto global (FMI)}
    {\footnotesize Según el FMI, los subsidios fósiles globales alcanzaron \textbf{USD 7 billones
    (7.1\% del PIB mundial)} en 2022, y los subsidios \emph{explícitos} más que se duplicaron
    \textbf{desde 2020} (Black, Liu, Parry \& Vernon, 2023).}
  \end{block}

  \note{%
    Este slide fija la lógica de identificación antes de cualquier dato. La clave de un experimento
    natural: la fuente de variación (el precio mundial) es ajena a la política que estudio. Insisto
    en describir vs.\ medir porque es el error típico: una serie que sube no prueba un efecto;
    necesito el
    contrafactual. El dato global del FMI lo atribuyo explícitamente a ellos, como contexto, no como
    hallazgo mío; y ojo: es ``desde 2020'', no ``en 2022''.}
\end{frame}

% ############################################################
% BLOQUE 2 — LOS DATOS
% ############################################################
\divisor{2}{Fuentes y cobertura del panel}{Los datos}

% --- Slide 3 (existente): Datos ---
\begin{frame}{Datos: panel de 34 países, 2015--2023}
  \small
  Construyo un panel anual de los países de América Latina y el Caribe con cobertura completa,
  combinando tres fuentes que cumplen funciones distintas:
  \begin{itemize}
    \item \textbf{Variable dependiente} (lo que quiero explicar): el subsidio explícito a
          combustibles fósiles, medido como \% del PIB. Viene de la \emph{IMF Fossil Fuel
          Subsidies Database}, la fuente obligatoria, que además separa el componente explícito
          del implícito.
    \item \textbf{Tratamiento} (la fuente de variación): el choque del precio del Brent en 2022
          (EIA), que entra al modelo \textbf{interactuado} con la condición de exportador neto.
          El choque por sí solo no identifica nada; lo hace su interacción con el grupo.
    \item \textbf{Variables fiscales} (deuda pública bruta, balance fiscal e ingreso del
          gobierno general, del \textbf{FMI, World Economic Outlook} vía IMF DataMapper): se usan
          \textbf{después} del modelo, en la parte de política. No entran a estimar el efecto
          porque serían ``malos controles'' (están afectadas por el propio choque).
  \end{itemize}

  \begin{block}{Estructura del panel}
    \key{34 países $\times$ 9 años $=$ 306 observaciones}, balanceado, sin huecos.
    \textbf{7 exportadores netos} frente a \textbf{27 importadores netos}.\\[0.3em]
    \textbf{Exportadores netos:} Bolivia, Colombia, Ecuador, Guyana, México, Trinidad y Tobago,
    Venezuela.
  \end{block}

  \note{%
    La base del FMI es la fuente obligatoria y distingue explícito de implícito; me centro en el
    explícito. Las fiscales las traje del World Economic Outlook del FMI (vía IMF DataMapper), no
    del Banco Mundial, porque cubre los 34 países sin huecos. Punto metodológico: las fiscales NO
    entran al modelo del efecto --serían
    malos controles--. Son solo 7 exportadores: grupo chico, con consecuencias en la precisión
    que reconozco abiertamente.}
\end{frame}

% --- Slide 2b (NUEVO): por qué World Economic Outlook sobre Banco Mundial + limitación de tamaño ---
\begin{frame}{Decisiones de fuente y sus límites}
  \small
  \begin{block}{Por qué el World Economic Outlook (FMI) y no el Banco Mundial}
    Las variables fiscales de la parte de política exigen \textbf{cobertura completa} de los 34
    países. El \emph{World Economic Outlook} del FMI la tiene; el Banco Mundial deja
    \alert{50--79\% de huecos en el Caribe}. Sin cobertura completa no se puede usar la deuda
    cuantitativamente (cruzarla contra el subsidio país por país), solo de forma anecdótica.\\[0.3em]
    {\footnotesize Indicadores WEO (vía IMF DataMapper), todos \% del PIB: deuda pública bruta
    (\texttt{GGXWDG\_NGDP}), balance fiscal (\texttt{GGXCNL\_NGDP}) e ingreso del gobierno general
    (\texttt{GGR\_G01\_GDP\_PT}).}
  \end{block}

  \begin{block}{Por qué el período 2015--2023}
    Centra la ventana alrededor del choque (pre-período estable 2015--2021, choque 2022,
    reversión 2023) y \textbf{excluye el inicio de la base}, donde la cobertura es irregular.
  \end{block}

  \textbf{Límite que reconozco de entrada:} hay solo \key{7 exportadores netos}. Es un grupo
  chico. Eso afecta la \textbf{precisión} (intervalos más anchos), \alert{no la insesgadez} (el
  estimador sigue centrado en el efecto verdadero). Lo retomo al discutir la significancia y la
  robustez.

  \note{%
    Aquí justifico dos decisiones de datos que un evaluador preguntaría. Primero, por qué el World
    Economic Outlook del FMI (vía IMF DataMapper) y no el Banco Mundial: no es preferencia, es
    cobertura. El Caribe tiene muchos huecos en el Banco
    Mundial (50--79\%), y sin la serie completa no puedo cruzar deuda contra subsidio para todos los
    países, que es justo lo que pide la parte fiscal. Segundo, la ventana 2015--2023 centra el
    choque y evita el arranque ruidoso de la base. Y dejo claro el límite del grupo chico: me cuesta
    precisión, no insesgadez. Esa distinción reaparece en los supuestos.}
\end{frame}

% --- Slide 2c (NUEVO): composición del panel — Tabla 2 ---
\begin{frame}{Quiénes componen el panel: los 34 países}
  \leavevmode
  {\small La muestra completa, clasificada por exposición petrolera neta y con su subsidio
  explícito año a año. \textbf{Qué mirar:} el \key{grupo de 7 exportadores} (arriba) frente a los
  27 importadores, y la heterogeneidad de niveles dentro de cada grupo (de Venezuela, con subsidio
  altísimo, a países con casi cero). El tratamiento se define por esta clasificación, no por
  ``país productor''.\par}\vspace{0.2em}
  \begin{center}
  \captitulo{Tabla 2}{Clasificación de los 34 países por exposición petrolera neta y subsidio explícito anual}
  \includegraphics[width=0.90\textwidth,height=0.60\textheight,keepaspectratio]{\tab{tab2_paises.pdf}}
  \end{center}

  \note{%
    Muestro la composición real del panel para que el grupo de tratamiento no sea abstracto. Lo que
    quiero que se vea: son solo 7 exportadores, y dentro de cada grupo hay enorme dispersión de
    niveles. Eso conecta con dos cosas dichas antes: la clasificación es por posición neta (Argentina
    y Brasil, aunque producen, son importadores netos de derivados), y el grupo chico de exportadores
    limita la precisión. Es referencia: no la leo entera, la dejo para quien quiera ver país por país.}
\end{frame}

% ############################################################
% BLOQUE 3 — LAS VARIABLES (y la justificación del método)
% ############################################################
\divisor{3}{Qué medimos y por qué}{Las variables y la estrategia de inferencia}

% --- Slide 2 (existente): Dos canales opuestos ---
\begin{frame}{Dos canales opuestos según la exposición petrolera}
  \small
  Un mismo choque de precios golpea en direcciones contrarias según de qué lado del mercado
  esté el país. Por eso la \textbf{posición petrolera neta} es la variable que organiza todo
  el análisis.
  \begin{itemize}
    \item \textbf{Exportador neto.} La renta petrolera financia buena parte del presupuesto, así
          que cuando el Brent sube, los ingresos fiscales se inflan. Ese ingreso extra abre
          espacio (y presión política) para \textbf{mantener los precios internos bajos} en
          lugar de trasladar el alza al consumidor. Resultado: el subsidio explícito se ensancha,
          y el país \textbf{sí tiene con qué pagarlo}.
    \item \textbf{Importador neto.} Compra el petróleo afuera, de modo que el alza encarece
          directamente su \textbf{costo de suministro}. Si el gobierno congela el precio al
          consumidor, el subsidio también crece, pero esta vez \textbf{sin renta extra} que lo
          respalde: el mismo choque que beneficia al exportador deteriora las cuentas del importador.
  \end{itemize}
  \textbf{Por qué ``neto'' y no ``productor'':} un país puede extraer crudo y aun así importar
  derivados; lo que decide quién gana o pierde renta es la posición \emph{neta}, no la producción.

  \begin{block}{Por qué el subsidio explícito (y no el implícito)}
    El \key{subsidio explícito} es la brecha entre el precio al consumidor y el costo de
    suministro: es precio-sensible y reacciona al choque. El \textbf{implícito} (externalidades
    ambientales más impuestos no cobrados) es estructural y no se mueve en el corto plazo
    $\Rightarrow$ lo uso como \textbf{prueba de placebo}: el efecto debe aparecer en el explícito
    y no en el implícito.
  \end{block}

  \note{%
    Aquí justifico la variable de tratamiento. Clasifico por exportador/importador \emph{neto}
    y no por ``país productor'' porque lo que importa para el efecto fiscal es la posición neta:
    un país puede producir petróleo y aun así ser importador neto de derivados. El explícito es
    precio-sensible (una brecha de precios); el implícito es estructural. Si el efecto aparece en
    el explícito y NO en el implícito, confirma que capturo el canal del precio: por eso lo uso
    como placebo más adelante.}
\end{frame}

% --- Slide 3c (NUEVO): por qué DiD, como cadena lógica ---
\begin{frame}{Por qué diferencias en diferencias (y no otra cosa)}
  \small
  No elegimos el método de un menú: el problema, leído con los datos, \textbf{conduce a un DiD}.
  La cadena:
  \begin{enumerate}
    \item La pregunta pide un \textbf{efecto} $\Rightarrow$ necesito un \textbf{contrafactual}
          (qué habría pasado sin el choque).
    \item El simple \textbf{antes/después} confunde el choque con la \emph{época}: en 2022 pasaron
          muchas cosas además del Brent.
    \item El simple \textbf{corte transversal} (exportadores vs.\ importadores en 2022) confunde el
          efecto con \emph{niveles preexistentes}: ya antes diferían.
    \item Las descriptivas revelan \key{dos dimensiones: tiempo $\times$ grupo}. Un grupo (los
          importadores) vivió la \textbf{misma época sin el canal} de la renta petrolera.
    \item Ese grupo es el \textbf{contrafactual}. Restar la ``época'' (lo común a 2022) al cambio
          de los expuestos \alert{es diferencias en diferencias}.
  \end{enumerate}

  \note{%
    Este es el slide bisagra entre describir y modelar. La idea que quiero transmitir: el DiD no es
    una elección de catálogo, es la consecuencia de tomarse en serio la pregunta. Antes/después solo
    confunde con la época; transversal solo confunde con niveles previos; necesito las dos
    diferencias a la vez. Y el grupo de control aparece naturalmente: los importadores vivieron 2022
    sin el canal de la renta. Restar la época a los expuestos es, literalmente, la doble diferencia.
    Así el método queda motivado por el problema, no impuesto.}
\end{frame}

% ############################################################
% BLOQUE 4 — ESTADÍSTICA DESCRIPTIVA
% ############################################################
\divisor{4}{Qué muestran los datos antes del modelo}{Estadística descriptiva}

% --- Slide 6 (existente): Ruptura 2022 — Figura 1 ---
\begin{frame}{El choque se sintió: ruptura clara en 2022}
  \begin{columns}[T]
    \begin{column}{0.46\textwidth}
      \small
      La figura traza el subsidio explícito de la región año a año, separado por
      grupo. Conviene leerla en dos pasos.\\[0.4em]
      \textbf{El agregado se quiebra en 2022:} de \key{0.85\% a 1.33\% del PIB}
      (2021$\to$2022), tras casi una década estable. En dólares, el explícito total
      de ALC pasó a $\sim$79 mil mill.\ (el implícito, estructural, $\sim$282 mil mill.).\\[0.4em]
      \textbf{Pero el salto no es parejo:} al desagregar por grupo, los exportadores
      saltan de \textbf{1.40\% a 2.54\%} (+1.14 pp) y los importadores apenas de
      \textbf{0.55\% a 0.66\%} (+0.12 pp). El alza de los exportadores fue
      \alert{$\sim$10$\times$} la de los importadores.\\[0.4em]
      Esa brecha visible es la que el modelo formaliza y cuantifica.
    \end{column}
    \begin{column}{0.54\textwidth}
      \centering
      \captitulo{Figura 1}{Subsidio explícito por grupo y año, 2015--2023}
      \includegraphics[width=\linewidth,height=0.60\textheight,keepaspectratio]{\fig{fig1_ruptura.png}}
      \fignota{\textit{Cómo leerla:} cada panel es la suma anual del subsidio explícito;
      (a)--(b) en USD miles de millones, (c)--(d) en \% del PIB. Fuente: IMF FFS Database.}
    \end{column}
  \end{columns}

  \note{%
    Empiezo por lo descriptivo porque cuenta la historia antes de cualquier modelo. La serie
    venía estable y en 2022 se quiebra. Lo importante no es el agregado: al separar por grupo, los
    exportadores se disparan y los importadores apenas se mueven. Aquí agrego ponderando por PIB
    (Tabla 1); el modelo trata cada país como una unidad, por eso los niveles del DiD son distintos
    pero la historia es la misma. Subrayo: la descriptiva es el diagnóstico que dictó el método, no
    decoración.}
\end{frame}

% --- Slide 6b (NUEVO): la ruptura en números — Tabla 1 ---
\begin{frame}{La ruptura, en números: subsidio por grupo y año}
  \leavevmode\par\vspace{-0.4em}
  {\footnotesize \textbf{Qué mirar:} la fila \emph{Explícito (\% PIB)}, columna \textbf{2022}
  frente a 2021, en cada panel. El total salta de 0.85 a \key{1.33}; el explícito de exportadores
  de 1.40 a \textbf{2.54} y el de importadores apenas de 0.55 a \textbf{0.66}. El \emph{implícito}
  no responde al choque (efecto $-0.32$, no significativo): sirve de placebo.\par}\vspace{0.1em}
  \begin{center}
  \captitulo{Tabla 1}{Subsidio explícito, implícito y total por grupo (USD y \% del PIB), 2015--2023}
  \includegraphics[width=0.92\textwidth,height=0.66\textheight,keepaspectratio]{\tab{tab1_descriptiva.pdf}}
  \end{center}

  \note{%
    Este slide da los números exactos detrás de la figura de ruptura, para quien quiera leerlos. La
    lectura guiada: fijarse en la fila del explícito en \% del PIB, columna 2022, en los tres
    paneles. El salto del total (0.85 a 1.33) y, sobre todo, el contraste exportadores (1.40 a 2.54)
    vs.\ importadores (0.55 a 0.66). Y noto que el efecto estimado sobre el implícito es nulo
    ($-0.32$, no significativo): aunque su serie varía en niveles, no responde al choque en el
    modelo, lo que anticipa su uso como placebo. La tabla agrega ponderando por PIB; el modelo
    trata cada país como unidad.}
\end{frame}

% --- Slide 7 (existente): Co-movimiento Brent — Figura 2 ---
\begin{frame}{Subsidio y precio del petróleo se mueven juntos}
  \begin{itemize}\small
    \item El Brent (izquierda) y el subsidio explícito agregado (derecha) trazan casi el
          mismo perfil, con \key{correlación de 0.75}: suben juntos hasta 2022 y bajan juntos en 2023.
    \item \textbf{Mecanismo:} cuando sube el costo de suministro y \textbf{no se traslada}
          al consumidor, la brecha (el subsidio) se ensancha de forma casi mecánica.
    \item \textbf{Pero es una pista, no una prueba:} el agregado \textbf{mezcla} a quien
          gana renta (exportadores) y a quien sufre el costo (importadores), que se mueven
          en sentidos opuestos. \alert{Motiva la pregunta causal; no la responde.}
  \end{itemize}
  \vspace{-0.6em}
  \begin{center}
    \captitulo{Figura 2}{Co-movimiento del Brent y el subsidio explícito agregado (doble eje), 2015--2023}
    \includegraphics[height=0.35\textheight]{\fig{fig2_brent.png}}
    \fignotaC{\textit{Cómo leerla:} doble eje, Brent (azul, izq.) y subsidio explícito de
    ALC (naranja, der.), 2015--2023. Importa el co-movimiento, no el nivel.
    Fuente: IMF FFS; EIA.}
  \end{center}

  \note{%
    Esta figura conecta tratamiento (Brent) y resultado (subsidio) de forma visual y honesta. La
    correlación de 0.75 es fuerte, pero correlación agregada $\neq$ causa: en el mismo número están
    sumados dos grupos con mecanismos opuestos. La uso para motivar, no para concluir. Marca la
    diferencia entre describir e identificar, que es lo que hace la parte siguiente. Las
    descriptivas son el diagnóstico que dictó el método.}
\end{frame}

% --- Slide 8b (NUEVO): heterogeneidad país a país — Figura 3 ---
\begin{frame}{Detrás del promedio: quién carga el efecto}
  \begin{columns}[T]
    \begin{column}{0.50\textwidth}
      \small
      El agregado esconde heterogeneidad. Ordenando los 34 países por cuánto cambió su
      subsidio con el choque, salta el patrón:
      \begin{itemize}
        \item \textbf{Los que más subieron:} \key{Surinam (+5.62), Bolivia (+5.47),
              Venezuela (+5.25)}. Casi todos los exportadores (naranja) van a la derecha.
        \item \textbf{No es exclusivo de exportadores:} Surinam y Argentina, \emph{importadores}
              (azul), están entre los más expuestos.
        \item \textbf{Hay contraejemplos:} Guyana, exportador, \alert{redujo} su subsidio
              ($-0.23$). El promedio del grupo no se cumple país por país.
      \end{itemize}
      \vspace{0.1em}
      Por eso el modelo trata cada país como una unidad, y la política cruza el efecto con el
      espacio fiscal de cada economía (Bloque 7), sin receta uniforme.
    \end{column}
    \begin{column}{0.50\textwidth}
      \centering
      \captitulo{Figura 3}{Cambio del subsidio por país, 2022 vs.\ 2015--2019}
      \includegraphics[width=\linewidth,height=0.66\textheight,keepaspectratio]{\fig{fig3_impacto.png}}
      \fignota{\textit{Cómo leerla:} cambio del subsidio explícito por país (pp del PIB),
      2022 vs.\ 2015--2019; naranja $=$ exportador, azul $=$ importador. Fuente: IMF FFS.}
    \end{column}
  \end{columns}

  \note{%
    Cierro la descriptiva mostrando la heterogeneidad que el promedio oculta. Tres lecturas: los que
    más suben son exportadores (Surinam, Bolivia, Venezuela), pero no es monopolio de exportadores
    (Surinam y Argentina importadores también están arriba), y hay contraejemplos como Guyana, que
    redujo su subsidio. Esto justifica dos decisiones: en el modelo trato cada país como unidad (no
    pondero por PIB como en el agregado), y en política cruzo el efecto con el espacio fiscal de cada
    economía en lugar de una recomendación uniforme. La figura es el puente entre la descriptiva y la
    parte de política.}
\end{frame}

% ############################################################
% BLOQUE 5 — EL MODELO
% ############################################################
\divisor{5}{Diferencias en diferencias: especificación y efecto}{El modelo}

% --- Slide 4 (existente): Estrategia DiD ---
\begin{frame}{Estrategia: diferencias en diferencias (DiD)}
  \begin{itemize}
    \item Comparo el \textbf{cambio} del subsidio en exportadores netos contra el \textbf{cambio}
          en importadores netos, del período pre-choque (2015--2021) al año del choque (2022).
  \end{itemize}

  \begin{block}{Especificación TWFE}
    \[
      \text{Subsidio}_{it} = \textcolor{AcentoRojo}{\boldsymbol{\beta_3}}\,(\text{Post2022}_t \times \text{Exportador}_i)
      + \alpha_i + \lambda_t + \varepsilon_{it}
    \]
    \begin{itemize}\small
      \item $\alpha_i$ (EF de país): absorben lo invariante del país, \textbf{incluido ser exportador}.
      \item $\lambda_t$ (EF de año): absorben lo común a cada año, \textbf{incluido el precio del Brent}.
      \item $\beta_3$ es el estimando: el cambio diferencial del subsidio de los exportadores atribuible al choque.
    \end{itemize}
  \end{block}

  \begin{itemize}\small
    \item \textbf{Qué mide la interacción:} $\text{Post2022}_t$ marca el momento del \textbf{choque
          de precios} y $\text{Exportador}_i$ identifica a los países cuya \textbf{posición petrolera
          neta} los expone a la renta. Su producto recoge la respuesta del subsidio al choque entre
          las economías que captan esa renta.
    \item Marco canónico 2$\times$2 (Roth et al., 2023). La identificación viene \textbf{solo de la
          interacción}: el nivel del Brent lo absorbe el EF de año. \textbf{Pre-registro:} la
          especificación se fijó \emph{antes} de estimar.
  \end{itemize}

  \note{%
    El truco del DiD: no basta ver que el subsidio de los exportadores subió --pudo subir por mil
    razones--. Lo comparo contra un control (importadores) que vivió el mismo año sin el canal de
    la renta petrolera. La doble diferencia limpia los factores comunes. Punto central: el Brent
    es el mismo para todos en cada año $\Rightarrow$ varianza cero dentro de año $\Rightarrow$ lo
    absorbe $\lambda_t$ y no lo identifico solo. La identificación viene EXCLUSIVAMENTE de la
    interacción. Si preguntan qué representa la interacción: Post2022 fecha el choque de precios y
    Exportador identifica la posición petrolera neta, que es la que expone al país a la renta. El
    producto mide entonces cómo responde el subsidio al choque entre las economías que captan esa
    renta. La clasificación es por posición neta, de modo que un país puede producir crudo y aun así
    quedar en el grupo de control si importa más derivados de los que exporta. Subrayo el
    pre-registro: definí esto antes de mirar resultados.}
\end{frame}

% --- Slide 9 (existente): Efecto central — Tabla 4 ---
\begin{frame}{El efecto central: $\sim$1.8 pp del PIB en los exportadores}
  \begin{columns}[T]
    \begin{column}{0.50\textwidth}
      \footnotesize
      \textbf{El número sale de cuatro promedios} (DiD crudo):
      \begin{itemize}
        \item Exportadores: +2.25 pp (3.30\,$\to$\,5.55)
        \item Importadores: +0.45 pp (0.43\,$\to$\,0.88)
        \item Doble diferencia $=$ \key{+1.80 pp del PIB}
      \end{itemize}
      \vspace{0.2em}
      La columna (2) de la tabla \textbf{confirma} ese crudo con efectos fijos:
      $\beta_3 = $ +1.79 pp, casi idéntico. Marginal al 10\% (SE 0.92, p $\approx$ 0.06):
      magnitud grande, precisión limitada por los 7 exportadores.
      \vspace{0.2em}

      \textbf{El placebo lo confirma: descarta un artefacto contable (columnas 3--4):}
      \begin{itemize}
        \item Explícito: +1.79 (responde al precio).
        \item Implícito: $-0.32$ (p = 0.52), \alert{nulo}.
      \end{itemize}
      Un sesgo contable afectaría \emph{ambos} componentes; que solo se mueva el
      explícito descarta esa explicación y confirma el canal del precio.
    \end{column}
    \begin{column}{0.50\textwidth}
      \centering
      \captitulo{Tabla 4}{Efecto del choque (DiD/TWFE): explícito, placebo y total}
      \includegraphics[width=\linewidth,height=0.52\textheight,keepaspectratio]{\tab{tab4_modelo.pdf}}
      \fignota{\textit{Cómo leerla:} cada columna es una especificación; la fila superior
      es $\beta_3$ con su EE entre paréntesis.}
    \end{column}
  \end{columns}

  \note{%
    Muestro primero el DiD ``a mano'' --las cuatro medias-- para que se vea de dónde sale el número
    antes de la caja negra del modelo: exportadores +2.25, importadores +0.45, diferencia 1.80. El
    modelo con EF da prácticamente lo mismo (1.79): el resultado no depende de la maquinaria. Sobre
    la significancia, directo: p $\approx$ 0.06, marginal al 10\%; con 7 exportadores no espero más.
    El golpe final es el placebo: efecto en el explícito, cero en el implícito. Si fuera un
    artefacto aparecería en ambos.}
\end{frame}

% ############################################################
% BLOQUE 6 — PRUEBAS DE LOS SUPUESTOS
% ############################################################
\divisor{6}{¿Se sostiene la identificación?}{Pruebas de los supuestos}

% --- Slide 5a (existente): Supuestos (1) tendencias paralelas ---
\begin{frame}{Supuestos de identificación (1): tendencias paralelas}
  \footnotesize
  \textbf{Tendencias paralelas:} \emph{sin el choque, ambos grupos habrían evolucionado igual.}
  \[
    \mathbb{E}[Y^0_{i,2022}-Y^0_{i,2021}\mid \text{Exp}{=}1]
    = \mathbb{E}[Y^0_{i,2022}-Y^0_{i,2021}\mid \text{Exp}{=}0]
  \]
  \begin{block}{Veredicto: EXAMINADO, se sostiene en lo sustantivo}
    \begin{itemize}
      \item Coeficientes pre-choque (2015--2020) alternan de signo
            ($+0.04,\,-0.41,\,-0.09,\,+0.95,\,+0.26,\,-0.77$) y todos sus \textbf{IC 95\% incluyen el cero}:
            no hay tendencia divergente que escale hacia 2022.
      \item \textbf{Matiz honesto:} el test conjunto de Wald \alert{rechaza} la nulidad estricta
            (\key{F = 4.09, p < 0.001}). Pero el rechazo viene de la \textbf{volatilidad} del grupo
            (7 exportadores; shocks 2018 y COVID-2020), \textbf{no de una tendencia} sistemática.
    \end{itemize}
  \end{block}
  \centering\itshape\color{GrisMedio}
  La volatilidad infla la incertidumbre; la tendencia sesgaría, y tendencia no hay.

  \note{%
    Esto es un examen, no una declaración de victoria. Lo que importa para que el DiD no esté
    sesgado es que no haya una \emph{tendencia} divergente previa, y no la hay: los coeficientes
    alternan de signo y sus intervalos cruzan el cero. El Wald rechaza que sean exactamente cero
    (p < 0.001), pero distingo dos cosas: el rechazo es por volatilidad --vaivenes de un grupo
    chico golpeado por 2018 y la pandemia--, no por una pendiente hacia 2022. La volatilidad me
    cuesta precisión, no insesgadez.}
\end{frame}

% --- Slide 8 (existente): Event study — Figura 4 (evidencia del supuesto 1) ---
\begin{frame}{La comparación es creíble: sin tendencia divergente previa}
  \begin{itemize}\small
    \item Cada punto es la brecha exportador--importador en un año, con 2021 fijado en cero.
          \textbf{A la izquierda de 2022 los puntos oscilan en torno al cero}: los dos grupos
          no venían separándose.
    \item \textbf{El salto ocurre justo en} \key{2022 (coef.\ +1.79)}: un quiebre, no la
          prolongación de una tendencia, \textbf{y se revierte en 2023} (+0.14, casi cero).
    \item El test conjunto rechaza la nulidad estricta (F = 4.09, p < 0.001), pero por
          \textbf{volatilidad} de un grupo chico, \textbf{no por una pendiente} hacia 2022.
  \end{itemize}
  \vspace{-0.6em}
  \begin{center}
    \captitulo{Figura 4}{Event study: efecto diferencial año a año (base 2021)}
    \includegraphics[height=0.38\textheight]{\fig{fig4_eventstudy.png}}
    \fignotaC{\textit{Cómo leerla:} efecto diferencial (pp del PIB), año a año, base 2021;
    barras $=$ IC 95\%. Fuente: IMF FFS Database.}
  \end{center}

  \note{%
    Aquí está la honestidad metodológica. Distinción crucial: una tendencia sistemática pre-choque
    (coeficientes que escalan) sesgaría el estimador; volatilidad alrededor de cero solo infla la
    incertidumbre. El test rechaza por lo segundo, no por lo primero: lo demuestro mostrando que
    los signos alternan. Argumento decisivo: la brecha no venía creciendo, brota en 2022 y
    desaparece en 2023 cuando el Brent baja. Por eso 2023 sale del estimador estático (es reversión)
    pero lo conservo en el event study para mostrar que el efecto se deshace.}
\end{frame}

% --- Slide 5b (existente): Supuestos (2) no anticipación ---
\begin{frame}{Supuestos de identificación (2): no anticipación}
  \textbf{No anticipación:} \emph{el subsidio de 2021 no reacciona a un choque de 2022 aún inexistente.}

  \begin{block}{Veredicto: NO testeable, defendido por diseño}
    \begin{itemize}
      \item No hay test estándar en un 2$\times$2.
      \item Argumento institucional: la \key{invasión de Ucrania (feb.\ 2022) fue inesperada};
            ningún gobierno reposicionó subsidios en 2021 anticipando el alza.
      \item La \textbf{exogeneidad del choque} es lo que sostiene este supuesto.
    \end{itemize}
  \end{block}

  \begin{center}
  \color{GrisMedio}\small
  El Brent es común a todos en cada año $\Rightarrow$ lo absorbe el EF de año ($\lambda_t$).\\
  La identificación viene \textbf{solo de la interacción} choque $\times$ exportador.
  \end{center}

  \note{%
    Este supuesto no se testea en un 2$\times$2, así que no finjo que lo validé con datos. Lo
    defiendo con el diseño: la guerra fue un shock inesperado; es implausible que alguien ajustara
    subsidios en 2021 por algo que no sabía que iba a pasar. Esa es la fortaleza de usar un choque
    exógeno. Cierro recordando que el Brent solo entra vía la interacción, porque su nivel lo come
    el efecto fijo de año.}
\end{frame}

% --- Slide 10 (existente): Robustez (leave-one-out) — Tabla 6 ---
\begin{frame}{El efecto es robusto: no lo carga un solo país}
  \begin{columns}[T]
    \begin{column}{0.52\textwidth}
      \small
      \textbf{Preocupación legítima:} ¿lo arrastra Venezuela, con un subsidio de 18.1\% del
      PIB? La tabla re-estima $\beta_3$ quitando los casos extremos, columna por columna:
      \begin{itemize}\footnotesize
        \item Muestra completa: 1.79$\dagger$
        \item Sin Venezuela: 1.15 (baja, pero \textbf{sigue positivo}).
        \item Sin Surinam: 1.98$^{*}$ (\emph{sube}): era un importador atípico que
              contaminaba el control.
        \item Sin ambos: 1.33$\dagger$
      \end{itemize}
      \vspace{0.2em}
      Y el \textbf{leave-one-out} (quitar cada exportador, uno a uno) mantiene $\beta_3$
      siempre positivo, en el rango \key{[1.15, 2.21]}.\\[0.3em]
      \textbf{Lectura:} Venezuela \alert{amplifica} el efecto, no lo crea. El resultado es
      robusto en signo y orden de magnitud, no un artefacto de un solo país.
    \end{column}
    \begin{column}{0.48\textwidth}
      \centering
      \captitulo{Tabla 6}{Robustez: $\beta_3$ excluyendo casos extremos}
      \includegraphics[width=\linewidth,height=0.48\textheight,keepaspectratio]{\tab{tab6_robustez.pdf}}
      \fignota{\textit{Cómo leerla:} cada columna es la misma especificación TWFE sobre una
      submuestra distinta. Si $\beta_3$ no cambia de signo, el efecto no depende de ese país.}
    \end{column}
  \end{columns}

  \note{%
    Anticipo la crítica: con un outlier como Venezuela, cualquiera sospecha que el resultado es
    solo Venezuela. En vez de quitarla a dedo --selección sobre el resultado, un error-- hago el
    leave-one-out de los siete. El efecto nunca cambia de signo y queda entre 1.15 y 2.21. Venezuela
    lo amplifica pero no lo inventa. Hallazgo interesante: quitar a Surinam \emph{sube} el efecto,
    porque es un importador con subsidio altísimo que contaminaba el control. Mantener la muestra
    completa pre-registrada como principal y mostrar la sensibilidad aparte es lo correcto.}
\end{frame}

% ############################################################
% BLOQUE 7 — RECOMENDACIONES DE POLÍTICA PÚBLICA
% ############################################################
\divisor{7}{Del coeficiente a la recomendación fiscal}{Recomendaciones de política pública}

% --- Slide 11 (existente): Matriz de política fiscal — Figura 5 ---
\begin{frame}{Del coeficiente al espacio fiscal: la matriz de política}
  \begin{columns}[T]
    \begin{column}{0.42\textwidth}
      \footnotesize
      Cada país se ubica por \textbf{cuánto pesa su subsidio} (eje horizontal) y
      \textbf{cuánta deuda} carga (eje vertical, su margen para sostenerlo). Las medianas
      (subsidio 0.81\%, deuda 64.1\%) parten el plano en cuatro cuadrantes.\\[0.4em]
      \textbf{El cuadrante crítico es el superior derecho:} subsidio alto \emph{y} deuda alta:
      \begin{itemize}\scriptsize
        \item \key{Reforma urgente:} Venezuela (18.1\,/\,164\,/\,$-5.3$), Surinam
              (8.5\,/\,112\,/\,$-2.6$), Bolivia (8.0\,/\,69\,/\,$-6.1$), Argentina.
        \item \textbf{Reforma gradual} (subsidio caro, deuda baja): Ecuador, Trinidad y Tobago,
              que tienen margen para una transición ordenada.
      \end{itemize}
      \vspace{0.2em}
      El tamaño del punto es el costo en dólares: $\sim$26 mil mill.\ en México, $\sim$6.8 en Colombia.
    \end{column}
    \begin{column}{0.58\textwidth}
      \centering
      \captitulo{Figura 5}{Matriz subsidio--deuda y cuadrantes de política (2022)}
      \includegraphics[width=\linewidth,height=0.55\textheight,keepaspectratio]{\fig{fig5_matriz_fiscal.png}}
      \fignota{\textit{Cómo leerla:} subsidio (eje X) vs.\ deuda (eje Y), \% del PIB 2022;
      tamaño $=$ costo del choque en USD. Arriba-derecha $=$ máxima presión fiscal.
      Fuente: IMF Fossil Fuel Subsidies Database; FMI, World Economic Outlook (vía IMF DataMapper).}
    \end{column}
  \end{columns}

  \note{%
    Esta es la pata de política, lo que el reto pide al final. ¿Por qué deuda y no balance? Porque
    el espacio fiscal en la literatura FMI/BID se ancla en el stock de deuda, no en el déficit de un
    año. El cuadrante crítico tiene nombre: Venezuela, Surinam, Bolivia, Argentina. Dos aclaraciones
    honestas: el costo en dólares es el efecto \emph{medio} aplicado a cada economía (un contrafactual,
    no el costo observado), y para Guyana, que redujo su subsidio, no reporto ese costo porque
    contradiría su dato. La presión fiscal NO es exclusiva de exportadores: Surinam y Argentina,
    importadores, están entre los más expuestos.}
\end{frame}

% --- Slide 11a (NUEVO): cómo se construye la matriz / los cuadrantes ---
\begin{frame}{Cómo se construye la matriz: los cuatro cuadrantes}
  \small
  La matriz es una herramienta de \textbf{priorización descriptiva}, no un modelo causal. Se arma en tres pasos:
  \begin{enumerate}
    \item \textbf{Dos ejes.} Cada país se ubica por su \textbf{subsidio explícito} (2022, \% PIB,
          eje horizontal) y su \textbf{deuda pública bruta} (2022, \% PIB, eje vertical).
    \item \textbf{Dos cortes.} Las \key{medianas de la muestra} (subsidio $=$ 0.81\%, deuda $=$
          64.1\%) dividen el plano. Uso medianas, no medias, porque son robustas a los extremos
          (Venezuela no arrastra el corte).
    \item \textbf{Cuatro cuadrantes} según en qué lado de cada mediana cae el país:
  \end{enumerate}
  \vspace{-0.2em}
  \begin{columns}[T]
    \begin{column}{0.49\textwidth}
      \begin{block}{Subsidio ALTO}
        \footnotesize
        \textbf{+ deuda alta $\Rightarrow$} \alert{Reforma urgente}\\
        \emph{(caro y sin margen para sostenerlo)}\\[0.3em]
        \textbf{+ deuda baja $\Rightarrow$} \alert{Reforma gradual}\\
        \emph{(caro, pero hay margen fiscal)}
      \end{block}
    \end{column}
    \begin{column}{0.49\textwidth}
      \begin{block}{Subsidio BAJO}
        \footnotesize
        \textbf{+ deuda alta $\Rightarrow$} \alert{Vigilar}\\
        \emph{(poco gasto, pero poco colchón)}\\[0.3em]
        \textbf{+ deuda baja $\Rightarrow$} \alert{Sin presión}\\
        \emph{(ni gasto alto ni deuda alta)}
      \end{block}
    \end{column}
  \end{columns}
  \vspace{0.1em}
  {\footnotesize \textbf{¿Por qué deuda y no balance?} El espacio fiscal se ancla en el
  \textbf{stock de deuda} acumulada, no en el déficit de un solo año, que es volátil.}

  \note{%
    Antes de mostrar la tabla, explico cómo se arma la matriz para que el cuadrante no parezca
    arbitrario. Tres pasos: dos ejes (subsidio y deuda, ambos \% PIB 2022), dos cortes en las
    medianas (0.81 y 64.1; uso medianas porque son robustas a outliers como Venezuela), y los cuatro
    cuadrantes según el lado de cada mediana. Insisto en que es priorización descriptiva, no causal: te
    dice dónde mirar primero, no qué política exacta aplicar. Y justifico el eje de deuda: el espacio
    fiscal es un tema de stock acumulado, no del déficit de un año. Esto prepara la lectura de la
    Tabla 5.}
\end{frame}

% --- Slide 11b (NUEVO): la matriz en detalle — Tabla 5 ---
\begin{frame}{La matriz, país por país: subsidio, deuda y cuadrante}
  \leavevmode
  {\small Aplicando la regla del slide anterior (medianas 0.81\% / 64.1\%), cada país queda en un
  cuadrante. \textbf{Qué mirar:} la columna \emph{Cuadrante}. El bloque de \key{``Reforma urgente''}
  (subsidio alto $+$ deuda alta) incluye a exportadores (Venezuela, Bolivia) \emph{y} a importadores
  (Surinam, Argentina): la presión no respeta la clasificación petrolera. El \emph{Costo del choque}
  es n.a.\ donde el efecto medio no aplica (Guyana lo redujo).
  {\color{GrisMedio}Fuente: IMF Fossil Fuel Subsidies Database; FMI, World Economic
  Outlook (vía IMF DataMapper).}\par}\vspace{0.2em}
  \begin{center}
  \captitulo{Tabla 5}{Espacio fiscal por país: subsidio, deuda, balance y cuadrante asignado (2022)}
  \includegraphics[width=0.95\textwidth,height=0.55\textheight,keepaspectratio]{\tab{tab5_fiscal.pdf}}
  \end{center}

  \note{%
    Esta tabla convierte la matriz visual en algo accionable: cada país con su subsidio, su deuda,
    su balance y el cuadrante que le toca. La lectura: la columna de cuadrante es la recomendación.
    Lo importante que la figura no muestra tan claro es el balance fiscal (varios de ``reforma
    urgente'' ya están en déficit fuerte: Bolivia $-6.1$, Venezuela $-5.3$). Y refuerzo el mensaje:
    en ``reforma urgente'' hay exportadores e importadores mezclados, así que la recomendación no se
    organiza por grupo petrolero sino por espacio fiscal. El costo en dólares es contrafactual
    (efecto medio aplicado), por eso es n.a.\ en Guyana, que fue a contracorriente.}
\end{frame}

% --- Slide 7b (NUEVO): regresividad / por qué focalizar ---
\begin{frame}{Por qué la reforma debe compensar de forma focalizada}
  \small
  Quitar o reducir un subsidio sube el precio del combustible y golpea a los hogares. La pregunta
  de política no es \emph{si} reformar, sino \textbf{cómo proteger a los hogares de menor ingreso}.

  \begin{block}{El subsidio actual es regresivo (FMI)}
    Según el FMI, el subsidio a los combustibles beneficia más a quien más consume, que son los
    hogares de mayor ingreso: el \key{20\% más rico capta más del 40\%} de los beneficios,
    \textbf{$\sim$6 veces} la porción del 20\% más pobre (Coady, Flamini \& Sears, 2015).
  \end{block}

  \textbf{Implicación de diseño:} un subsidio universal al precio es un instrumento mal dirigido.
  Sustituirlo por \alert{transferencias focalizadas} a los hogares vulnerables \textbf{mejora la
  equidad} (llega a quien lo necesita) \textbf{y} libera espacio fiscal (deja de subsidiar al
  decil más alto). La reforma no tiene por qué ser regresiva si la compensación está bien dirigida.

  \note{%
    Este slide sostiene la recomendación de focalizar. El dato de regresividad lo atribuyo
    explícitamente al FMI (Coady et al.\ 2015), no como hallazgo mío: el 20\% más rico se lleva más
    del 40\%, unas 6 veces lo del 20\% más pobre. La lección de diseño: el subsidio universal al
    precio es plata mal dirigida; la transferencia focalizada protege al vulnerable y ahorra al
    fisco a la vez. Eso desactiva la objeción de que ``reformar castiga al pobre'': castiga al pobre
    si se hace mal; bien diseñada, es progresiva.}
\end{frame}

% --- Slide 11c (NUEVO): hoja de ruta por cuadrante ---
\begin{frame}{Hoja de ruta: qué hacer en cada cuadrante}
  \small
  La recomendación \textbf{se diferencia por cuadrante}, no es uniforme. Una secuencia general en
  tres pasos, calibrada según el margen fiscal:
  \begin{block}{Acción por cuadrante}
    \footnotesize
    \begin{itemize}
      \item \alert{Reforma urgente} (Venezuela, Surinam, Bolivia, Argentina): ajuste
            \textbf{prioritario} del precio interno con \textbf{compensación focalizada inmediata}.
            El déficit ya es fuerte (\key{Bolivia $-6.1$, Venezuela $-5.3$}) $\Rightarrow$ no hay
            margen para diferir.
      \item \alert{Reforma gradual} (Ecuador, Trinidad y Tobago): hay deuda baja $\Rightarrow$
            \textbf{secuencia escalonada} con cronograma anunciado y compensación que suavice el alza.
      \item \alert{Vigilar / Sin presión:} mantener monitoreo. El riesgo es que un nuevo choque de
            precios los empuje al cuadrante crítico (el subsidio es \textbf{pro-cíclico}).
    \end{itemize}
  \end{block}
  \vspace{-0.2em}
  {\footnotesize \textbf{Secuencia transversal:} (1) congelar la indexación del subsidio al precio
  internacional; (2) liberar el precio de forma gradual y anunciada; (3) redirigir lo ahorrado a
  \textbf{transferencias focalizadas} a los hogares vulnerables (no universales, por la regresividad).}

  \note{%
    Este es el slide que convierte el coeficiente y la matriz en acción concreta, que es lo que el
    reto pide. La idea central: la recomendación se diferencia por cuadrante. En reforma urgente no
    hay margen (los déficit ya son fuertes, Bolivia y Venezuela), hay que actuar ya con compensación
    inmediata. En reforma gradual hay colchón de deuda baja, así que se puede secuenciar con
    cronograma. En los otros dos, vigilar, porque el subsidio es pro-cíclico y un nuevo choque los
    mueve al cuadrante malo. Y doy una secuencia operativa de tres pasos: congelar indexación,
    liberar precio gradual, redirigir el ahorro a transferencias focalizadas. La focalización (no
    universal) viene de la regresividad del slide anterior.}
\end{frame}

% --- Slide 12 (existente): Conclusión y recomendación ---
\begin{frame}{Conclusión y recomendación de política}
  \begin{itemize}
    \item \textbf{Hallazgo:} el choque de 2022 encareció \textbf{diferencialmente} el subsidio
          explícito de los exportadores netos en \key{$\sim$1.8 pp del PIB}: efecto robusto,
          concentrado en el componente sensible al precio, y \textbf{transitorio} (se revierte en 2023).
    \item \textbf{Matiz:} la vulnerabilidad fiscal del subsidio \textbf{precede al choque} y alcanza
          a importadores (Surinam, Argentina). El choque \textbf{agravó} una exposición preexistente.
  \end{itemize}

  \begin{block}{Recomendación: reforma diferenciada por espacio fiscal (no uniforme)}
    \begin{itemize}
      \item Donde coinciden subsidio caro y deuda alta $\Rightarrow$ \alert{reforma prioritaria}.
      \item Donde hay margen fiscal $\Rightarrow$ \textbf{transición gradual con compensación
            focalizada} a hogares vulnerables.
    \end{itemize}
    \vspace{0.2em}
    {\footnotesize El subsidio al combustible es \textbf{regresivo} (lo capta más quien más
    consume, es decir los hogares de mayor ingreso), así que sustituirlo por transferencias
    focalizadas mejora la equidad además del balance fiscal.}
  \end{block}

  \begin{center}\color{Carbon}\small
  Un subsidio que se infla con cada choque de precios es un \textbf{pasivo fiscal pro-cíclico}:
  reformarlo libera espacio fiscal y reduce la exposición a la volatilidad del petróleo.
  \end{center}

  \note{%
    Cierro uniendo evidencia y política. El hallazgo en una frase: el choque golpea más fuerte a los
    exportadores, pero de forma transitoria. El matiz evita la lectura simplista: no es solo problema
    de exportadores; el subsidio caro ya era riesgo fiscal en varios importadores. La recomendación
    no es ``eliminar subsidios'' en abstracto, sino diferenciar por espacio fiscal. Mensaje que quiero
    que quede: el subsidio es pro-cíclico, sube justo cuando el petróleo sube, que es cuando más cuesta
    sostenerlo; reformarlo es ganar resiliencia fiscal. Ese es el gancho de política tipo BID.}
\end{frame}

% ============================================================
% REFERENCIAS
% ============================================================
\begin{frame}{Referencias}
  \footnotesize
  \begin{itemize}
    \item Black, S., Liu, A. A., Parry, I., \& Vernon, N. (2023). \emph{IMF Fossil Fuel Subsidies
          Data: 2023 Update}. IMF Working Paper No.\ 2023/169.
    \item Coady, D., Flamini, V., \& Sears, L. (2015). \emph{The Unequal Benefits of Fuel Subsidies
          Revisited}. IMF Working Paper WP/15/250.
    \item Roth, J., Sant'Anna, P., Bilinski, A., \& Poet, J. (2023). What's trending in
          difference-in-differences? A synthesis of the recent econometrics literature.
          \emph{Journal of Econometrics}.
    \item International Monetary Fund. \emph{Fossil Fuel Subsidies Database} (2023 update).
          \textbf{[Fuente de datos principal]}
    \item International Monetary Fund. \emph{World Economic Outlook Database} (acceso vía IMF
          DataMapper API). Indicadores: deuda pública bruta (\texttt{GGXWDG\_NGDP}), resultado
          fiscal (\texttt{GGXCNL\_NGDP}) e ingreso del gobierno general (\texttt{GGR\_G01\_GDP\_PT}),
          todos como \% del PIB. \textbf{[Variables fiscales]}
    \item U.S. Energy Information Administration. \emph{Europe Brent Spot Price}.
          \textbf{[Precio del petróleo]}
  \end{itemize}

  \note{%
    No leo la lista; es respaldo. Las anclas son tres: el contexto del choque y la regresividad
    (FMI: Black et al.\ 2023 y Coady et al.\ 2015), el marco DiD y sus supuestos (Roth et al.\ 2023),
    y las fuentes de datos (IMF FFS Database y EIA). Mantengo la bibliografía corta a propósito: es
    una prueba técnica, no una tesis.}
\end{frame}

% ============================================================
% (El antiguo "Material de respaldo" se eliminó: sus 4 piezas
%  — tab1, tab2, fig3, tab5 — se integraron al cuerpo principal
%  con lectura guiada en sus bloques respectivos.)
% ============================================================

% ============================================================
% RESPALDO — PREGUNTAS ANTICIPADAS (Q&A de defensa)
% Solo se usan si surge la pregunta en la sustentación.
% ============================================================
{
\setbeamertemplate{footline}{} % divisor sin footline
\begin{frame}[plain]
  \vfill
  \begin{center}
    {\color{AcentoRojo}\rule{2.6cm}{1.2pt}}\\[0.9em]
    {\Large\color{Carbon}\bfseries Preguntas anticipadas}\\[0.5em]
    {\color{GrisMedio}Material de respaldo para la defensa}\\[0.9em]
    {\color{AcentoRojo}\rule{2.6cm}{1.2pt}}
  \end{center}
  \vfill
\end{frame}
}

% --- Q1: Wald rechaza tendencias paralelas ---
\begin{frame}{P1. ``El Wald rechaza tendencias paralelas, ¿por qué confías en el DiD?''}
  \small
  \begin{block}{Hay que distinguir dos violaciones distintas}
    \begin{itemize}
      \item \textbf{Tendencia divergente sistemática} (coeficientes que \emph{escalan} hacia 2022)
            $\Rightarrow$ \alert{sesga} el estimador. Es la grave.
      \item \textbf{Volatilidad alrededor de cero} (vaivenes que alternan de signo) $\Rightarrow$
            infla la incertidumbre, \alert{no sesga}.
    \end{itemize}
  \end{block}
  \begin{itemize}
    \item El Wald (\key{F = 4.09, p < 0.001}) rechaza la nulidad \emph{estricta} (que los
          coeficientes pre-choque sean exactamente cero), no detecta una tendencia. La forma del
          rechazo lo confirma: los coeficientes alternan de signo ($+0.04,-0.41,-0.09,+0.95,+0.26,-0.77$),
          \textbf{no escalan}.
    \item El ruido viene de un grupo de solo 7 países (shocks 2018 y COVID-2020): cuesta
          \textbf{precisión} (por eso p $\approx$ 0.06 en el efecto), no insesgadez.
    \item \textbf{Decisivo:} la brecha no venía creciendo, \textbf{brota en 2022 y desaparece en
          2023}. Una tendencia preexistente no hace eso; un quiebre ligado al choque, sí.
  \end{itemize}

  \note{%
    La respuesta en un minuto: el Wald rechaza, sí, pero lo importante es \emph{qué} rechaza. Rechaza
    que los coeficientes pre-choque sean exactamente cero, no que haya una tendencia. Y la diferencia
    es todo: una tendencia que escala sesga; volatilidad que alterna solo resta precisión. Lo pruebo
    con la forma del rechazo —los signos alternan, no escalan— y con el patrón temporal: brota en
    2022 y se va en 2023. Eso es un quiebre, no una divergencia previa. Con 7 países el ruido es
    inevitable; me cuesta el p de 0.06, no la validez.}
\end{frame}

% --- Q2: tratamiento continuo vs binario ---
\begin{frame}{P2. ``¿Y si usas tratamiento continuo en vez de binario exportador?''}
  \small
  \begin{block}{Qué aportaría}
    Un tratamiento continuo (exposición petrolera neta como \% del PIB, interactuada con Post2022)
    capturaría la \textbf{intensidad}, no solo el signo. Predicción: efecto monótono en la dosis
    (los más expuestos reaccionan más) $\Rightarrow$ sería evidencia \textbf{dosis-respuesta} que
    refuerza el mecanismo.
  \end{block}
  \textbf{Por qué elegí el binario (tres razones defendibles):}
  \begin{itemize}
    \item La hipótesis es sobre el \textbf{signo} del canal (exportador gana renta / importador no):
          una distinción cualitativa.
    \item Con solo \key{7 exportadores}, el continuo gasta grados de libertad que no tengo: el
          binario es más robusto con $N$ chico.
    \item \textbf{Honestidad:} no estimé la versión continua, así que no afirmo su resultado.
  \end{itemize}
  \textbf{Matiz técnico:} el \emph{choque} (Brent) es el mismo para todos; lo que varía es la
  \emph{exposición} estructural. Eso se modela mejor como heterogeneidad del efecto (binario
  interactuado) que como ``dosis'' (por eso descarté un DiD de tratamiento continuo a la Callaway).
  \alert{Es una extensión futura, no un resultado obtenido.}

  \note{%
    Aquí soy honesto sobre un límite. El continuo daría dosis-respuesta, que sería lindo. Pero lo
    descarté por tres razones: la hipótesis es de signo, no de intensidad; con 7 exportadores no me
    sobran grados de libertad; y no lo corrí, así que no voy a fingir que sé qué da. El matiz fino:
    el choque es común a todos, lo que varía es la exposición estructural, y eso se captura mejor
    como heterogeneidad del efecto que como dosis del tratamiento. Lo dejo planteado como extensión,
    no como hallazgo.}
\end{frame}

% --- Q3: reversión de 2023 ---
\begin{frame}{P3. ``El coeficiente de 2023 es 0.14, ¿el efecto se revirtió del todo?''}
  \small
  \begin{itemize}
    \item Sí, en lo esencial. El punto cae de \key{1.79 (2022) a 0.14 (2023)}: $\sim$92\% de reversión.
    \item El IC 95\% de 2023 es $[-0.65,\,+0.92]$: incluye holgadamente el cero y es estrecho en
          torno a él $\Rightarrow$ en 2023 el efecto \textbf{no se distingue de ``sin efecto''}.
          (Para contraste, el de 2022 es $[-0.01,\,+3.58]$.)
    \item \textbf{Matiz:} ``revertido'' significa que el \textbf{efecto diferencial} desapareció (los
          exportadores dejaron de subsidiar \emph{más} que los importadores), no que el subsidio
          volviera a su nivel exacto pre-choque.
  \end{itemize}
  \begin{block}{Por qué esto apoya el mecanismo}
    El Brent cayó de \textbf{101 a 82 en 2023 ($-18\%$)}: se apagó el motor del efecto. Que se
    revierta cuando el precio baja es evidencia \alert{a favor} del mecanismo causal, no en contra.
    Por eso 2023 sale del estimador estático pero se conserva en el event study, para mostrar la
    transitoriedad.
  \end{block}

  \note{%
    La respuesta corta: sí, ~92\% de reversión, y el intervalo de 2023 abraza el cero y es estrecho,
    así que ese año no hay efecto distinguible. El matiz que evita malentendidos: revertido es que el
    \emph{diferencial} se apagó, no que el subsidio volviera al céntimo de antes. Y el punto de fondo:
    se revierte justo cuando el Brent baja de 101 a 82. Que el efecto siga al precio es la mejor
    señal de que es el precio lo que lo causa. Por eso 2023 lo saco del estimador estático pero lo
    muestro en el event study.}
\end{frame}

% --- Q4: Bolivia régimen estructural ---
\begin{frame}{P4. ``Bolivia ya subsidiaba mucho antes, ¿cómo separas el choque del régimen?''}
  \small
  \begin{itemize}
    \item Los \textbf{efectos fijos de país} hacen exactamente eso: absorben todo lo constante de
          Bolivia (gas estatal, precios históricamente congelados, instituciones). El DiD \textbf{no
          compara el nivel} de Bolivia con nadie: compara su \emph{cambio} contra el de los importadores.
    \item Los datos: nivel pre-choque $\approx$ 2.64\% del PIB (ese es el régimen estructural,
          irrelevante para el DiD). Lo que el DiD usa es el salto de \key{+5.38 pp} en 2022 (de 2.64
          a 8.02\% del PIB).
    \item \textbf{¿Y si habría saltado igual sin el choque?} Ahí entra el contrafactual: si fuera
          estructural, lo veríamos crecer \emph{antes} de 2022 (no se ve: tendencias paralelas) y
          \textbf{no se revertiría}. Pero Bolivia vuelve a 5.03\% en 2023. \alert{Un régimen
          estructural no se revierte; un efecto de precio sí.}
  \end{itemize}
  \begin{block}{En una frase}
    El nivel estructural lo absorbe el efecto fijo de país; lo estimado es el \textbf{cambio
    diferencial} en el año del choque, y su reversión en 2023 confirma que es el precio, no el régimen.
  \end{block}

  \note{%
    La clave es que el DiD nunca mira niveles, mira cambios. Todo lo que hace a Bolivia un caso
    estructural —el gas estatal, los precios congelados— es constante en el tiempo, así que lo come
    el efecto fijo de país. Lo que estimo es el salto de +5.38 puntos en 2022, no el nivel de 2.64.
    Y si alguien insiste en que ese salto era estructural y no del choque, el contrafactual responde:
    un régimen estructural habría venido creciendo (no lo hace) y no se revertiría (Bolivia baja a
    5.03 en 2023). La reversión es la prueba: el precio prende y apaga el efecto; un régimen, no.}
\end{frame}

\end{document}
